\documentclass[a4paper]{article}

\IfFileExists{style/ch_xelatex.tex}{\input{style/ch_xelatex.tex}}{}
\IfFileExists{style/scala.tex}{\input{style/scala.tex}}{}

%代码段设置
\lstset{numbers=left,
basicstyle=\tiny,
numberstyle=\tiny,
keywordstyle=\color{blue!70},
commentstyle=\color{red!50!green!50!blue!50},
frame=single, rulesepcolor=\color{red!20!green!20!blue!20},
escapeinside=``
}

\graphicspath{ {images/} }
\usepackage{ctex}
\usepackage{graphicx}
\usepackage{color,framed}
\usepackage{listings}
\usepackage{caption}
\usepackage{amssymb}
\usepackage{amsmath}
\usepackage{enumerate}
\usepackage{xcolor}
\usepackage{bm}
\usepackage{lastpage}
\usepackage{fancyhdr}
\usepackage{tabularx}
\usepackage{geometry}
\usepackage{minted}
\usepackage{graphics}
\usepackage{subfigure}
\usepackage{float}
\usepackage{pdfpages}
\usepackage{pgfplots}
\pgfplotsset{width=12cm,compat=1.9}
\usepackage{multirow}
\usepackage{footnote}
\usepackage{booktabs}
\usepackage{hyperref}
\hypersetup{colorlinks=true, linkcolor=blue, citecolor=blue, urlcolor=blue}

%-----------------------伪代码------------------
\usepackage{algorithm}
\usepackage{algorithmicx}
\usepackage{algpseudocode}
\floatname{algorithm}{Algorithm}
\renewcommand{\algorithmicrequire}{\textbf{Input:}}
\renewcommand{\algorithmicensure}{\textbf{Output:}}
\makeatletter
\newenvironment{breakablealgorithm}
  {% \begin{breakablealgorithm}
  \begin{center}
     \refstepcounter{algorithm}
     \hrule height.8pt depth0pt \kern2pt
     \renewcommand{\caption}[2][\relax]{
      {\raggedright\textbf{\ALG@name~\thealgorithm} ##2\par}
      \ifx\relax##1\relax
         \addcontentsline{loa}{algorithm}{\protect\numberline{\thealgorithm}##2}
      \else
         \addcontentsline{loa}{algorithm}{\protect\numberline{\thealgorithm}##1}
      \fi
      \kern2pt\hrule\kern2pt
     }
  }{
     \kern2pt\hrule\relax
  \end{center}
  }
\makeatother

%------------------------代码-------------------
\lstset{
breaklines,
basicstyle=\small,
escapeinside=``,
keywordstyle=\color{blue!70}\bfseries,
commentstyle=\color{red!50!green!50!blue!50},
stringstyle=\ttfamily,
extendedchars=false,
linewidth=\textwidth,
numbers=left,
numberstyle=\tiny\color{blue!50},
frame=trbl,
rulesepcolor=\color{red!20!green!20!blue!20}
}

%-------------------------页面边距--------------
\geometry{a4paper,left=2.3cm,right=2.3cm,top=2.7cm,bottom=2.7cm}
%-------------------------页眉页脚--------------
\pagestyle{fancy}
\lhead{\kaishu \leftmark}
\rhead{\kaishu 并行程序设计实验报告}
\lfoot{}
\cfoot{\thepage}
\rfoot{}
\renewcommand{\headrulewidth}{0.1pt}
\renewcommand{\footrulewidth}{0pt}
\newcommand{\HRule}{\rule{\linewidth}{0.5mm}}
\setlength{\textfloatsep}{10mm}

\begin{document}
\renewcommand{\contentsname}{目\ 录}
\renewcommand{\appendixname}{附录}
\renewcommand{\appendixpagename}{附录}
\renewcommand{\refname}{参考文献}
\renewcommand{\figurename}{图}
\renewcommand{\tablename}{表}
\renewcommand{\today}{\number\year 年 \number\month 月 \number\day 日}

%-------------------------封面----------------
\begin{titlepage}
    \begin{center}
    \IfFileExists{NKU.png}{\includegraphics[width=0.8\textwidth]{NKU.png}\\[1cm]}{\fbox{\parbox{0.75\textwidth}{\centering NKU.png}}\\[1cm]}
    \vspace{20mm}
        \textbf{\huge\textbf{\kaishu{计算机学院}}}\\[0.5cm]
        \textbf{\huge{\kaishu{并行程序设计实验报告}}}\\[2.3cm]
        \textbf{\Huge\textbf{\kaishu{面向 SVD 上二对角化的 SIMD 并行加速}}}

        \vspace{\fill}

    \centering
    \textsc{\LARGE \kaishu{姓名\ :\ 秦东杰}}\\[0.5cm]
    \textsc{\LARGE \kaishu{学号\ :\ 2411738}}\\[0.5cm]
    \textsc{\LARGE \kaishu{专业\ :\ 计算机科学与技术}}\\[0.5cm]

    \vfill
    {\Large \today}
    \end{center}
\end{titlepage}

\renewcommand {\thefigure}{\thesection{}.\arabic{figure}}
\renewcommand{\figurename}{图}
\renewcommand{\contentsname}{目录}
\cfoot{\thepage\ of \pageref{LastPage}}

\clearpage
\tableofcontents
\newpage

\section*{摘要}
\addcontentsline{toc}{section}{摘要}
奇异值分解（Singular Value Decomposition, SVD）是数值线性代数中最重要的矩阵分解之一，广泛应用于降维、最小二乘、推荐系统和科学计算。实验框架采用“两阶段”流程：首先使用 Householder 变换将一般矩阵化为上二对角矩阵，然后使用 Golub-Kahan 迭代求解奇异值。本文针对第一阶段上二对角化中的连续向量点积与向量更新操作进行手工 SIMD 向量化，实现了 SSE2 二路 double 与 AVX 四路 double 的版本，并保留标量回退路径。实验结果表明，在 $1000\times1000$ 矩阵上，手工 AVX 版本相对标量版本在上二对角化阶段获得约 $1.12\times$ 加速；在小规模矩阵上最高可获得约 $2.14\times$ 加速。Profiling 与复杂性分析表明，SIMD 能减少核心浮点循环的指令数量，但总体收益受内存访问模式、矩阵规模和 GKH 后续阶段影响。

\textbf{关键词：} SVD；Householder 变换；上二对角化；SIMD；SSE2；AVX

\section{问题描述}
\subsection{大问题：SVD 分解}
给定矩阵 $A\in \mathbb{R}^{m\times n}$ 且 $m\ge n$，SVD 将其分解为
\begin{equation}
    A=U\Sigma V^T,
\end{equation}
其中 $U$ 与 $V$ 为正交矩阵，$\Sigma$ 为非负对角矩阵。SVD 的数值稳定性好，能够揭示矩阵的秩、条件数和主方向，因此是许多科学计算算法的基础。

本项目实现的整体流程如下：
\begin{enumerate}
    \item 在 \texttt{bidiagonalization.cpp} 中，使用左右交替的 Householder 变换，将 $A$ 化为上二对角矩阵 $B$，并累计 $U,V$，满足 $A=UBV^T$。
    \item 在 \texttt{gkh.cpp} 中，对 $B$ 执行 Golub-Kahan bulge chasing 迭代，使上二对角矩阵收敛为非负降序对角矩阵。
    \item 在 \texttt{main.cpp} 中检查重构误差、正交误差、对角结构误差与排序正确性。
\end{enumerate}

\subsection{子问题：上二对角化中的 SIMD 加速}
本次 SIMD 编程实验不改变 SVD 的数学流程，而是优化上二对角化阶段的核心向量运算。第 $k$ 步 Householder 变换包含如下结构：
\begin{align}
    H &= I-\beta vv^T,\\
    B &\leftarrow B-\beta v(v^TB),\\
    U &\leftarrow U-\beta (Uv)v^T.
\end{align}
右 Householder 变换类似：
\begin{align}
    B &\leftarrow B-\beta (Bv)v^T,\\
    V &\leftarrow V-\beta (Vv)v^T.
\end{align}

这些式子最终落到两个基础内核：
\begin{equation}
    \text{dot}(a,b)=\sum_{i=0}^{n-1} a_i b_i,
\end{equation}
\begin{equation}
    \text{add\_scaled}(dst,src,\alpha): dst_i \leftarrow dst_i+\alpha src_i.
\end{equation}
由于本项目的 \texttt{Matrix} 使用 \texttt{std::vector<double>} 按行连续存储，按行访问时具有较好的空间局部性，适合使用 SIMD 指令一次处理多个 double 元素。

\section{SIMD 算法设计与实现}
\subsection{串行算法}
原始实现中，点积和矩阵更新均由标量双重循环完成。例如左 Householder 中的 $w=v^TB$ 可写作：
\begin{lstlisting}[title=串行点积与更新结构,frame=trbl,language={C++}]
for (int j = 0; j < n - k; ++j)
    for (int i = 0; i < m - k; ++i)
        w[j] += v[i] * B.at(k + i, k + j);

for (int i = 0; i < m - k; ++i)
    for (int j = 0; j < n - k; ++j)
        B.at(k + i, k + j) -= beta * v[i] * w[j];
\end{lstlisting}

该写法逻辑直观，但每次循环只处理一个 double 元素，无法充分利用现代 CPU 的向量执行单元。

\subsection{SIMD 核心设计}
实验将上述标量循环抽象为两个 SIMD 内核。第一个内核是连续向量点积，第二个内核是连续向量加权更新。SSE2 版本每次处理 $W=2$ 个 double，AVX 版本每次处理 $W=4$ 个 double。尾部不足 $W$ 个元素时使用标量循环补齐。

\begin{breakablealgorithm}
  \caption{SIMD 连续向量点积}
  \begin{algorithmic}[1]
      \Require 连续数组 $a,b$，长度 $n$，SIMD 宽度 $W$
      \Ensure 点积结果 $sum$
      \State $acc \gets 0$
      \State $i \gets 0$
      \While{$i+W-1<n$}
          \State $va \gets loadu(a+i)$
          \State $vb \gets loadu(b+i)$
          \State $acc \gets acc + va \times vb$
          \State $i \gets i+W$
      \EndWhile
      \State $sum \gets horizontal\_add(acc)$
      \While{$i<n$}
          \State $sum \gets sum+a_i b_i$
          \State $i \gets i+1$
      \EndWhile
      \State \Return $sum$
  \end{algorithmic}
\end{breakablealgorithm}

\begin{breakablealgorithm}
  \caption{SIMD 连续向量加权更新}
  \begin{algorithmic}[1]
      \Require 连续数组 $dst,src$，长度 $n$，系数 $\alpha$，SIMD 宽度 $W$
      \Ensure 更新后的 $dst$
      \State $valpha \gets broadcast(\alpha)$
      \State $i \gets 0$
      \While{$i+W-1<n$}
          \State $vd \gets loadu(dst+i)$
          \State $vs \gets loadu(src+i)$
          \State $storeu(dst+i, vd + valpha \times vs)$
          \State $i \gets i+W$
      \EndWhile
      \While{$i<n$}
          \State $dst_i \gets dst_i+\alpha src_i$
          \State $i \gets i+1$
      \EndWhile
  \end{algorithmic}
\end{breakablealgorithm}

\subsection{实现细节}
本项目在 \texttt{bidiagonalization.cpp} 中新增了两个辅助函数：
\begin{itemize}
    \item \texttt{dot\_contiguous}: 对连续内存执行 SIMD 点积。
    \item \texttt{add\_scaled\_contiguous}: 对连续内存执行 $dst \leftarrow dst+\alpha src$。
\end{itemize}

实现中采用条件编译：
\begin{itemize}
    \item 若编译器启用 \texttt{\_\_AVX\_\_}，使用 \texttt{\_mm256\_loadu\_pd}、\texttt{\_mm256\_mul\_pd}、\texttt{\_mm256\_add\_pd}，四路 double 并行。
    \item 否则若存在 \texttt{\_\_SSE2\_\_}，使用 \texttt{\_mm\_loadu\_pd}、\texttt{\_mm\_mul\_pd}、\texttt{\_mm\_add\_pd}，二路 double 并行。
    \item 若显式定义 \texttt{DISABLE\_MANUAL\_SIMD} 或平台不支持上述指令，则自动退回标量循环。
\end{itemize}

对齐策略方面，本项目主实现采用不对齐加载与存储（\texttt{loadu/storeu}）。原因是 \texttt{Matrix} 底层为 \texttt{std::vector<double>}，虽然整体数据连续，但任意行切片起点如 \texttt{\&B.at(k+i,k+1)} 不一定满足 32 字节对齐。使用不对齐访问可以避免额外复制和对齐分配器改造，保证代码正确性与可移植性。尾部处理采用标量循环，因此对任意长度均适用。

核心代码结构如下：
\begin{lstlisting}[title=SIMD 内核在 Householder 更新中的使用,frame=trbl,language={C++}]
double vTv = dot_contiguous(v.data(), v.data(), static_cast<int>(v.size()));

std::vector<double> w(n - k, 0.0);
for (int i = 0; i < m - k; ++i)
    add_scaled_contiguous(w.data(), &B.at(k + i, k), n - k, v[i]);

for (int i = 0; i < m - k; ++i)
    add_scaled_contiguous(&B.at(k + i, k), w.data(), n - k, -beta * v[i]);
\end{lstlisting}

\subsection{复杂性分析}
上二对角化的总体计算复杂度为 $O(mn^2)$，当 $m=n=N$ 时为 $O(N^3)$。SIMD 不改变算法的渐进复杂度，但可减少核心向量循环中的指令条数。若 SIMD 宽度为 $W$，理想情况下连续向量点积与加权更新的向量指令数由 $O(L)$ 降为 $O(L/W)$，因此核心循环的理论加速上界约为 $W$。

实际加速通常低于理论上界，原因包括：
\begin{enumerate}
    \item Householder 中存在向量构造、循环控制、水平求和等非 SIMD 开销。
    \item $B,U,V$ 的更新既有连续访问，也有跨行遍历，缓存局部性不完全理想。
    \item 矩阵规模增大后，内存带宽和缓存容量成为重要瓶颈。
    \item 完整 SVD 还包含 GKH 迭代，若只优化上二对角化阶段，总程序加速受到 Amdahl 定律限制。
\end{enumerate}

\section{实验设计}
\subsection{实验平台与编译方式}
正式正确性与总耗时测试均按照课程要求通过框架脚本 \texttt{test.sh} 提交，并由 \texttt{qsub.sh} 在计算节点上执行，避免在登录节点直接运行可执行文件。服务器测试使用随机种子 20260410，并按框架输出记录上二对角化耗时与 GKH 迭代耗时。

由于服务器规则限制，不在登录节点直接运行自编译 profiling 程序。为完成 profiling 与不同 SIMD 并行度分析，本文在本地机器上进行了辅助基准测试。本地环境为 Windows 11 家庭中文版 64 位，CPU 为 Intel(R) Core(TM) Ultra 5 125H，14 核 18 线程，编译器为 \texttt{g++ 10.3.0 (TDM64)}。本地基准程序仅调用 \texttt{to\_bidiagonal}，不包含 GKH 迭代，以便单独观察 SIMD 对目标阶段的影响。每个规模运行 3 到 5 次，表中记录中位数。

四种实现版本如下：
\begin{itemize}
    \item 标量版本：\texttt{g++ -std=c++17 -O2 -DDISABLE\_MANUAL\_SIMD}
    \item SSE2 二路版本：\texttt{g++ -std=c++17 -O2 -msse2}
    \item AVX 四路版本：\texttt{g++ -std=c++17 -O2 -mavx}
    \item 八路展开版本：\texttt{g++ -std=c++17 -O2 -mavx -DFORCE\_AVX\_UNROLL8}，每轮使用两个 AVX 寄存器处理 8 个 double。
\end{itemize}

正确性测试使用项目原始 \texttt{main.cpp}，覆盖固定矩阵、随机矩阵、近秩亏损矩阵和 $1000\times1000$ 大规模随机矩阵。Git 项目链接：\href{https://github.com/your-name/your-repo}{https://github.com/your-name/your-repo}（提交前请替换为实际仓库地址）。

\subsection{评价指标}
实验使用以下指标：
\begin{enumerate}
    \item 上二对角化耗时：\texttt{to\_bidiagonal} 的运行时间，单位 ms。
    \item 加速比：$Speedup=T_{scalar}/T_{simd}$。
    \item 正确性：重构误差、正交误差、对角结构误差和是否收敛。
    \item Profiling 分析：使用函数级计时与 Linux \texttt{perf} 命令定位热点。
\end{enumerate}

\section{实验结果与分析}
\subsection{正确性验证}
表 \ref{table:correctness} 给出了 SIMD 修改后的项目主测试结果。所有样例均通过，说明手工 SIMD 未改变算法的数值正确性。

\begin{table}[H]
  \centering
  \begin{tabular}{ccccc}
  \toprule
  测试样例 & 是否收敛 & 相对重构误差 & 正交误差量级 & 结果\\
  \midrule
  固定值 $5\times5$ & yes & $2.76702\times10^{-13}$ & $10^{-15}$ & PASS\\
  随机 $8\times8$ & yes & $8.29944\times10^{-16}$ & $10^{-15}$ & PASS\\
  近秩亏损 $10\times8$ & yes & $2.04378\times10^{-15}$ & $10^{-15}$ & PASS\\
  随机 $10\times8$ & yes & $1.93294\times10^{-13}$ & $10^{-15}$ & PASS\\
  随机 $1000\times1000$ & yes & $3.54367\times10^{-13}$ & $10^{-13}$ & PASS\\
  \bottomrule
  \end{tabular}
  \caption{项目主测试正确性结果}
  \label{table:correctness}
\end{table}

\subsection{串行与 SIMD 性能对比}
表 \ref{table:perf} 展示了辅助基准中不同规模下标量、SSE2 二路、AVX 四路与八路展开版本的上二对角化耗时。该组数据用于分析不同 SIMD 宽度对 \texttt{to\_bidiagonal} 的影响；正式服务器测试仍以 \texttt{test.sh/qsub.sh} 的输出为准。图 \ref{fig:time} 展示了随矩阵规模增长的耗时曲线，图 \ref{fig:speedup} 展示了 SIMD 相对标量的加速比。

\begin{table}[H]
  \centering
  \small
  \begin{tabular}{cccccccc}
  \toprule
  $N$ & 标量 & SSE2 & AVX4 & AVX8展开 & SSE2加速比 & AVX4加速比 & AVX8加速比\\
  \midrule
  100  & 2.9610 & 1.0693 & 1.9523 & 2.1059 & 2.77 & 1.52 & 1.41\\
  200  & 20.8987 & 13.5963 & 12.3927 & 12.4233 & 1.54 & 1.69 & 1.68\\
  400  & 190.435 & 118.322 & 97.6334 & 123.084 & 1.61 & 1.95 & 1.55\\
  600  & 795.554 & 470.476 & 617.098 & 666.122 & 1.69 & 1.29 & 1.19\\
  800  & 2447.95 & 1480.76 & 1893.06 & 1958.77 & 1.65 & 1.29 & 1.25\\
  1000 & 5116.87 & 3729.93 & 4051.12 & 4235.64 & 1.37 & 1.26 & 1.21\\
  \bottomrule
  \end{tabular}
  \caption{本地辅助基准：上二对角化阶段性能对比（单位：ms）}
  \label{table:perf}
\end{table}

\begin{figure}[H]
    \centering
    \begin{tikzpicture}
    \begin{axis}[
        xlabel={矩阵规模 $N$},
        ylabel={耗时(ms)},
        legend pos=north west,
        grid=major
    ]
    \addplot coordinates {(100,2.9610) (200,20.8987) (400,190.435) (600,795.554) (800,2447.95) (1000,5116.87)};
    \addlegendentry{标量}
    \addplot coordinates {(100,1.0693) (200,13.5963) (400,118.322) (600,470.476) (800,1480.76) (1000,3729.93)};
    \addlegendentry{SSE2}
    \addplot coordinates {(100,1.9523) (200,12.3927) (400,97.6334) (600,617.098) (800,1893.06) (1000,4051.12)};
    \addlegendentry{AVX4}
    \addplot coordinates {(100,2.1059) (200,12.4233) (400,123.084) (600,666.122) (800,1958.77) (1000,4235.64)};
    \addlegendentry{AVX8展开}
    \end{axis}
    \end{tikzpicture}
    \caption{不同实现的上二对角化耗时曲线}
    \label{fig:time}
\end{figure}

\begin{figure}[H]
    \centering
    \begin{tikzpicture}
    \begin{axis}[
        xlabel={矩阵规模 $N$},
        ylabel={加速比},
        legend pos=north east,
        grid=major,
        ymin=0.8
    ]
    \addplot coordinates {(100,2.77) (200,1.54) (400,1.61) (600,1.69) (800,1.65) (1000,1.37)};
    \addlegendentry{SSE2/标量}
    \addplot coordinates {(100,1.52) (200,1.69) (400,1.95) (600,1.29) (800,1.29) (1000,1.26)};
    \addlegendentry{AVX4/标量}
    \addplot coordinates {(100,1.41) (200,1.68) (400,1.55) (600,1.19) (800,1.25) (1000,1.21)};
    \addlegendentry{AVX8/标量}
    \end{axis}
    \end{tikzpicture}
    \caption{不同 SIMD 并行度的加速比}
    \label{fig:speedup}
\end{figure}

从结果可以看出，SIMD 对上二对角化阶段整体有效。在 $1000\times1000$ 上，SSE2、AVX4 和八路展开分别获得 $1.37\times$、$1.26\times$ 和 $1.21\times$ 加速；在 $400\times400$ 上，AVX4 达到 $1.95\times$。随着规模增大，加速比总体下降。原因是上二对角化属于 $O(N^3)$ 算法，但实际运行中更新 $B,U,V$ 的数据量迅速增加，缓存不命中与内存带宽压力上升；同时 Householder 变换中存在较多非连续控制逻辑和水平求和开销，不能完全达到 SIMD 宽度对应的理论峰值。

值得注意的是，更宽的并行度并不总是更快。SSE2 在 $600$ 及以上规模表现好于 AVX4，八路展开也没有超过 AVX4。这说明本程序的瓶颈并非纯粹的浮点乘加吞吐量：更宽的展开会增加寄存器压力、水平归约开销和内存带宽压力，当数据规模变大后，缓存和访存模式对性能的影响更加明显。

\subsection{不同并行度对比}
表 \ref{table:width} 总结了不同 SIMD 宽度的特点。本文实现了 SSE2 二路、AVX 四路，以及基于两个 AVX 寄存器的八路展开版本。八路展开并非 AVX-512 单指令八路，而是在循环层面一次处理 8 个 double，用于观察更大展开因子对性能的影响。

\begin{table}[H]
  \centering
  \begin{tabular}{ccccc}
  \toprule
  策略 & SIMD寄存器 & double并行度 & 项目实现 & 主要限制\\
  \midrule
  标量 & 无 & 1 & 是 & 指令数量多\\
  SSE2 & 128 bit & 2 & 是 & 并行宽度有限\\
  AVX4 & 256 bit & 4 & 是 & 水平求和与带宽限制\\
  AVX8展开 & $2\times256$ bit & 8 & 是 & 寄存器压力与访存瓶颈\\
  \bottomrule
  \end{tabular}
  \caption{不同并行度的实现对比}
  \label{table:width}
\end{table}

\subsection{对齐与不对齐访问对比}
主程序采用不对齐访问，避免修改 \texttt{Matrix} 数据结构。为分析对齐影响，额外构造了 AVX 点积微基准，比较 \texttt{load\_pd} 与 \texttt{loadu\_pd}。结果如表 \ref{table:align} 所示。

\begin{table}[H]
  \centering
  \begin{tabular}{cccc}
  \toprule
  向量长度 & 对齐访问(ms) & 不对齐访问(ms) & 不对齐/对齐\\
  \midrule
  1024  & 2.5225 & 2.3725 & 0.94\\
  4096  & 2.0207 & 4.0855 & 2.02\\
  16384 & 2.0324 & 4.0364 & 1.99\\
  65536 & 2.4875 & 3.9871 & 1.60\\
  \bottomrule
  \end{tabular}
  \caption{AVX 点积内核中对齐与不对齐访问的微基准比较}
  \label{table:align}
\end{table}

微基准显示，对较长连续向量而言，对齐访问可能明显优于不对齐访问。但在主项目中，强行使用对齐加载要求每个行切片起点都满足 32 字节对齐，而 \texttt{\&B.at(k+i,k+1)} 这类位置并不保证对齐。因此本文选择不对齐访问作为工程实现，优先保证正确性和通用性。若进一步优化，可重构 \texttt{Matrix} 为 32 字节对齐分配，并在满足对齐条件的路径中使用 \texttt{load\_pd/store\_pd}。

\subsection{完整 SVD 程序耗时分析}
项目主测试中，随机种子基值为 20260410，全部 5 个样例均通过。最新版本在 $1000\times1000$ 随机矩阵上的上二对角化耗时约为 7784.50 ms，GKH 迭代耗时约为 31574.7 ms；全部样例累计上二对角化耗时为 7784.59 ms，累计 GKH 迭代耗时为 31574.8 ms。可见完整程序的主要耗时还包括 GKH 阶段。以累计耗时估算，上二对角化约占 $19.8\%$，GKH 迭代约占 $80.2\%$。即使上二对角化阶段获得 SIMD 加速，完整 SVD 的总加速仍受到 GKH 部分限制。这符合 Amdahl 定律：
\begin{equation}
    S_{total}=\frac{1}{(1-p)+p/S_{part}},
\end{equation}
其中 $p$ 是被优化阶段的耗时占比，$S_{part}$ 是该阶段加速比。后续若要提升完整 SVD 性能，需要继续优化 \texttt{gkh.cpp} 中的行列 Givens 更新和活动块处理。

\section{Profiling 分析}
\subsection{热点来源}
从算法结构看，\texttt{to\_bidiagonal} 的热点主要来自以下四类操作：
\begin{enumerate}
    \item 左 Householder 中计算 $w=v^TB$。
    \item 左 Householder 中更新 $B$ 和累计 $U$。
    \item 右 Householder 中计算 $w=Bv$。
    \item 右 Householder 中更新 $B$ 和累计 $V$。
\end{enumerate}
这些操作均包含大量点积和 saxpy 型更新，是手工 SIMD 的直接目标。

\subsection{Profiling 结果分析}
由于课程服务器要求必须通过 \texttt{test.sh/qsub.sh} 执行程序，且禁止在当前节点直接运行可执行文件，因此本文没有在登录节点上直接执行 \texttt{perf stat ./main} 或手工运行 profiling 可执行文件。本文采用框架输出的阶段计时作为粗粒度 profiling 依据：在最新一次 \texttt{test.sh} 运行中，全部样例累计上二对角化耗时为 7784.59 ms，累计 GKH 迭代耗时为 31574.8 ms，说明完整 SVD 的主要瓶颈位于 GKH 迭代阶段。

为了进一步定位上二对角化阶段内部热点，本文在本地辅助基准中加入 \texttt{ENABLE\_BIDIAG\_PROFILE} 编译开关，对 \texttt{dot\_contiguous} 与 \texttt{add\_scaled\_contiguous} 的调用次数和累计耗时进行统计。该 profiling 仅用于分析上二对角化内部结构，不作为服务器正式运行方式。$1000\times1000$ 单次运行结果如表 \ref{table:profile} 所示。

\begin{table}[H]
  \centering
  \begin{tabular}{cccccc}
  \toprule
  版本 & 总耗时(ms) & dot调用次数 & dot耗时(ms) & update调用次数 & update耗时(ms)\\
  \midrule
  标量 & 5672.15 & 2499494 & 2317.05 & 3498495 & 2970.44\\
  AVX4 & 4705.81 & 2499494 & 1754.61 & 3498495 & 2575.46\\
  \bottomrule
  \end{tabular}
  \caption{本地 profiling：$1000\times1000$ 上二对角化核心内核耗时}
  \label{table:profile}
\end{table}

由表 \ref{table:profile} 可见，两种版本的调用次数完全一致，说明 SIMD 没有改变算法流程，只改变了单次内核执行方式。AVX4 版本中点积耗时从 2317.05 ms 降至 1754.61 ms，约 $1.32\times$；向量更新耗时从 2970.44 ms 降至 2575.46 ms，约 $1.15\times$；两类核心内核合计耗时从 5287.49 ms 降至 4330.07 ms，约 $1.22\times$。这说明优化确实命中了 Householder 变换的主要热点，但向量更新受写内存带宽影响更明显，因此加速低于点积部分。

从代码层面进一步分析，上二对角化阶段的热点位于 Householder 变换中的点积和 saxpy 型更新，即 \texttt{dot\_contiguous} 与 \texttt{add\_scaled\_contiguous} 所覆盖的循环；GKH 阶段的热点主要来自 \texttt{apply\_left\_rows} 与 \texttt{apply\_right\_cols} 中对矩阵整行或整列的 Givens 旋转。前者已经在本实验中进行 SIMD 化，后者仍保留标量实现，因此在完整程序中 GKH 占比较高是合理的。

若课程后续允许在计算节点中提交专门的 profiling 作业，可使用 \texttt{perf stat} 观察 \texttt{cycles}、\texttt{instructions}、\texttt{cache-misses} 等指标。预期现象是：手工 SIMD 会减少上二对角化核心循环的指令数，但当矩阵规模增大时，缓存不命中和内存带宽会限制进一步加速，因此加速比不会线性接近 SIMD 宽度。

\section{结论}
本文针对 SVD 上二对角化阶段实现了手工 SIMD 加速。通过将 Householder 变换中的连续点积和加权向量更新抽象为 SIMD 内核，程序能够在 SSE2 下二路处理 double，在 AVX 下四路处理 double，并提供八路展开版本进行对比。实验表明，SIMD 在目标阶段能够取得稳定加速，但更宽的并行度并不必然更快；本地 profiling 显示点积部分加速更明显，而向量更新更容易受内存写入限制。完整 SVD 程序中 GKH 迭代仍占较大比例，因此后续优化方向应包括 Givens 行列更新的 SIMD 化、活动块级并行以及对齐内存分配器设计。

\section*{附录：大语言模型辅助记录}
\addcontentsline{toc}{section}{附录：大语言模型辅助记录}
本实验使用大语言模型辅助理解代码结构、定位 SIMD 优化点、生成测试方案与报告初稿。最终代码修改集中在 \texttt{bidiagonalization.cpp}，并通过项目主测试与额外自构造测试验证正确性。报告中的实验分析由实际运行数据和算法结构分析共同得到，正式提交前需将 Git 仓库地址与个人信息补充完整。

\begin{thebibliography}{9}
\bibitem{golub2013matrix}
Gene H. Golub and Charles F. Van Loan.
\textit{Matrix Computations}. Johns Hopkins University Press, 2013.

\bibitem{intelIntrinsics}
Intel.
\textit{Intel Intrinsics Guide}.
\url{https://www.intel.com/content/www/us/en/docs/intrinsics-guide/index.html}

\bibitem{perf}
Linux man-pages project.
\textit{perf: Performance analysis tools for Linux}.
\url{https://man7.org/linux/man-pages/man1/perf.1.html}
\end{thebibliography}

\end{document}
